\renewcommand{\Anno}{2008}

\renewcommand{\NumeroProblema}{4}

\renewcommand{\Testo}{\begin{enumerate}
    \item Alcuni anni fa, dopo un mese di lezione, una piccola rivoluzione è avvenuta all'interno della scuola Sant'Anna: un po' di studenti è infatti passato dalla classe di Ingegneria a quella di Medicina.  I responsabili delle due classi sono comunque rimasti contenti in quanto a seguito dei trasferimenti il Quoziente Intellettivo (QI) medio degli studenti di entrambe le classi è aumentato.  Dopo un altro mese di lezione è avvenuta un'altra piccola rivoluzione: un po' di studenti sono questa volta passati da Medicina ad Ingegneria.  È possibile che il QI medio in ciascuna delle due classi sia aumentato ancora una volta?

    \item L'anno scorso si è verificato un fenomeno analogo, che questa volta ha coinvolto tre classi.  Dopo un mese un po' di studenti sono passati da Ingegneria a Medicina, ed un po' di studenti sono passati da Medicina a Giurisprudenza.  Ancora una volta il QI medio in ciascuna delle tre classi è aumentato. Dopo un altro mese è avvenuto uno spostamento in senso inverso: un po' di studenti sono passati da Giurisprudenza a Medicina, ed un po' di studenti sono passati da Medicina ad Ingegneria.  È possibile che ancora una volta il QI medio in ciascuna delle tre classi sia aumentato?
\end{enumerate}

{\small Nota: si assuma che il QI di una persona sia un intero positivo che
non varia nel tempo.}
}

\renewcommand{\Risposta}{
\begin{enumerate}
	\item[(a)]  Non è possibile.

	\item[(b)]  È possibile.
\end{enumerate}
}

\renewcommand{\Aiutino}{
\begin{enumerate}
	\renewcommand{\labelenumi}{(\alph{enumi})} 
	\item[(a)] Indichiamo con $I_{1}$ ed $I_{2}$ il QI medio degli studenti di Ingegneria prima e dopo il primo turno di trasferimenti.  Definiamo $M_{1}$ ed $M_{2}$ in maniera analoga, facendo riferimento agli studenti di Medicina.  Indichiamo con $T_{1}$ il QI medio degli studenti che al primo colpo si sono trasferiti da Ingegneria a Medicina, e con $T_{2}$ il QI medio degli studenti che si trasferiscono da Medicina ad Ingegneria al secondo colpo.
	
	Supponiamo per assurdo che entrambi i trasferimenti migliorino entrambe le classi.  Quali disuguaglianze sussistono tra $I_{1}$, $M_{1}$, $T_{1}$?  E tra $I_{2}$, $M_{2}$, $T_{2}$?  Che disuguaglianza sussiste tra $T_{1}$ ed $M_{2}$?

	\item[(b)] Basta esibire un esempio in cui la situazione descritta si realizza.  L'idea è più o meno la seguente.  Ad Ingegneria ci sono tanti studenti medi e qualcuno buono.  A Medicina ci sono tanti studenti scarsi, pochi buoni e pochi buonissimi.  Da Ingegneria un gran numero di studenti medi passa a Medicina, permettendo così ad uno buonissimo di andare a Giurisprudenza, dove sono tutti molto buoni.  A questo punto basta che un qualunque studente di Giurisprudenza passi a Medicina, ed uno studente buono di Medicina passi ad Ingegneria.
\end{enumerate}
}

\renewcommand{\SoluzioneUfficiale}{\TitoloSoluzione \DomandaA
Non è possibile.  Per dimostrarlo ragioniamo per assurdo. Indichiamo con $I_{1}$ ed $I_{2}$ il QI medio degli studenti di Ingegneria prima e dopo il primo turno di trasferimenti.  Definiamo $M_{1}$ ed $M_{2}$ in maniera analoga, facendo riferimento agli studenti di Medicina.  Indichiamo con $T_{1}$ il QI medio degli studenti che al primo colpo si sono trasferiti da Ingegneria a Medicina.

Dalle ipotesi sappiamo che
$$
    I_{1}<I_{2},
    \hspace{3em}
    M_{1}<M_{2}.
$$

È chiaro che questo può accadere solo se quelli che si sono trasferiti erano peggiori della media di Ingegneria e migliori della media di Medicina, cioè
$$
M_{1}<T_{1}<I_{1}.
$$

Inoltre la nuova media a Medicina sarà minore di $T_{1}$, perché come abbiamo appena visto i restanti iscritti a Medicina sono mediamente peggio di quelli che si sono appena trasferiti. In simboli
$$
    M_{2}<T_{1}.
$$

Supponiamo ora che anche la seconda ondata di trasferimenti migliori il QI medio di entrambe le classi. Indicando con $T_{2}$ il QI medio degli studenti che si trasferiscono, questo è possibile se e solo se
$$
    I_{2}<T_{2}<M_{2}.
$$

Mettendo insieme le disuguaglianze finora dimostrate si avrebbe però che
$$M_{2}<T_{1}<I_{1}<I_{2}<T_{2}<M_{2},$$
il che è chiaramente assurdo.

\DomandaB 
È possibile.  Per dimostrarlo basta esibire un esempio in cui la situazione descritta si realizza.  Consideriamo inizialmente 3 classi
con la seguente composizione\medskip

\hfill
{\renewcommand{\arraystretch}{1.3}
\begin{tabular}{|c|c|}
    \hline
    \multicolumn{2}{|c|}{\textbf{Ingegneria}}  \\
    \hline
    Studenti & QI  \\
    \hline\hline
    50 & 2  \\
    \hline
    1 & 3  \\
    \hline\hline
    \multicolumn{2}{|l|}{Totale: 51 studenti}  \\
    \multicolumn{2}{|l|}{QI medio: $\sim$ 2.02}  \\
    \hline
\end{tabular}
\hfill
\begin{tabular}{|c|c|}
    \hline
    \multicolumn{2}{|c|}{\textbf{Medicina}}  \\
    \hline
    Studenti & QI  \\
    \hline\hline
    100 & 1  \\
    \hline
    1 & 5  \\
    \hline
    1 & 10  \\
    \hline\hline
    \multicolumn{2}{|l|}{Totale: 102 studenti}  \\
    \multicolumn{2}{|l|}{QI medio: $\sim$ 1.13}  \\
    \hline
\end{tabular}
\hfill
\begin{tabular}{|c|c|}
    \hline
    \multicolumn{2}{|c|}{\textbf{Giurisprudenza}}  \\
    \hline
    Studenti & QI  \\
    \hline\hline
    1 & 9  \\
    \hline
    1 & 10  \\
    \hline\hline
    \multicolumn{2}{|l|}{Totale: 2 studenti}  \\
    \multicolumn{2}{|l|}{QI medio: 9.5}  \\
    \hline
\end{tabular}
}
\hfill\mbox{}
\medskip

Supponiamo che al primo turno 40 studenti con QI uguale a 2 passino da Ingegneria a Medicina, e lo studente con QI uguale a 10 passi da Medicina a Giurisprudenza.  In questo modo è chiaro che Ingegneria e Giurisprudenza ci guadagnano. Facendo i conti si vede che pure Medicina ci guadagna, in quanto la partenza dello studente con QI pari a 10 viene compensata dall'arrivo di ben 40 studenti con QI pari a 2.

Dopo la prima ondata di trasferimenti la situazione è dunque la
seguente:\medskip

\hfill
{\renewcommand{\arraystretch}{1.3}
\begin{tabular}{|c|c|}
    \hline
    \multicolumn{2}{|c|}{\textbf{Ingegneria}}  \\
    \hline
    Studenti & QI  \\
    \hline\hline
    10 & 2  \\
    \hline
    1 & 3  \\
    \hline\hline
    \multicolumn{2}{|l|}{Totale: 11 studenti}  \\
    \multicolumn{2}{|l|}{QI medio: $\sim$ 2.09}  \\
    \hline
\end{tabular}
\hfill
\begin{tabular}{|c|c|}
    \hline
    \multicolumn{2}{|c|}{\textbf{Medicina}}  \\
    \hline
    Studenti & QI  \\
    \hline\hline
    100 & 1  \\
    \hline
    40 & 2  \\
    \hline
    1 & 5  \\
    \hline\hline
    \multicolumn{2}{|l|}{Totale: 141 studenti}  \\
    \multicolumn{2}{|l|}{QI medio: $\sim$ 1.31}  \\
    \hline
\end{tabular}
\hfill
\begin{tabular}{|c|c|}
    \hline
    \multicolumn{2}{|c|}{\textbf{Giurisprudenza}}  \\
    \hline
    Studenti & QI  \\
    \hline\hline
    1 & 9  \\
    \hline
    2 & 10  \\
    \hline\hline
    \multicolumn{2}{|l|}{Totale: 3 studenti}  \\
    \multicolumn{2}{|l|}{QI medio: $\sim$ 9.67}  \\
    \hline
\end{tabular}
}
\hfill\mbox{}
\medskip

Supponiamo ora che nella seconda ondata di trasferimenti lo studente di Giurisprudenza con QI uguale a 9 passi a Medicina, e che lo studente di Medicina con QI uguale a 5 passi ad Ingegneria. È immediato verificare che tutte le tre classi ci guadagnano.

Se qualcuno ritiene che i QI di questo esempio siano troppo bassi per studenti della Scuola Sant'Anna, basta che aggiunga 100 o 150 a tutti, senza cambiare la situazione!
}

\renewcommand{\SoluzioniAlternative}{}

\renewcommand{\Osservazioni}{\TitoloOsservazione Non è difficile
dimostrare che la situazione descritta al punto (b) è impossibile se tutti i 4 trasferimenti coinvolgono ciascuno uno studente solo (o
comunque se tutti i 4 trasferimenti coinvolgono lo \emph{stesso}
numero di studenti). Con ovvio significato dei simboli si avrebbe infatti che
$$T_{I\to M}<I_{1}<I_{2}<T_{M\to I}<T_{G\to M}<
G_{2}<T_{M\to G}<T_{I\to M},$$
il che è ovviamente assurdo. Le singole disuguaglianze si dimostrano con ragionamenti analoghi a quelli utilizzati per il punto (a).
}

\renewcommand{\Commenti}{
}